% Phụ lục C

\chapter{Bảng phân công công việc} % Tên của phụ lục

\label{AppendixC} % Để trích dẫn chương này ở chỗ nào đó trong bài, hãy sử dụng lệnh \ref{AppendixC}

%----------------------------------------------------------------------------------------

\begin{table}[htbp]
	\centering
	\caption{Bảng phân công công việc và đóng góp của các thành viên Nhóm 6.}
	\label{tab:phancông_nhom6}
	\footnotesize % Giảm cỡ chữ xuống một chút để khít hoàn toàn vào trang giấy, không bị tràn viền
	\begin{tabularx}{\textwidth}{|c|p{2.8cm}|c|X|c|c|c|}
		\hline
		\textbf{STT} & \textbf{Họ và tên} & \textbf{Mã SV} & \textbf{Phân công nhiệm vụ} & \begin{tabular}[c]{@{}c@{}}\textbf{Mức độ}\\\textbf{hoàn thành}\end{tabular} & \begin{tabular}[c]{@{}c@{}}\textbf{Trọng số}\\\textbf{cá nhân}\end{tabular} & \textbf{Chữ ký} \\ \hline
		1 & Nguyễn Doãn Bình & 23001582 & Nội dung quyển báo cáo \newline Nội dung Slide & 80\% & 24.50\% & Bình \\ \hline
		2 & Phan Thành Đạt & 23001595 & Nội dung quyển báo cáo \newline Nội dung Slide \newline Thuyết trình & 80\% & 24.50\% & Đạt \\ \hline
		3 & Tô Quang Hân & 23001605 & Nội dung quyển báo cáo \newline Nội dung Slide \newline Code mô hình RandomForest và LinearSVC & 87\% & 25.50\% & Hân \\ \hline
		4 & Nguyễn Trọng Tùng & 23001638 & Nội dung quyển báo cáo \newline Giao diện và nội dung slide \newline Code mô hình Logistics, CNN, LSTM \newline Tinh chỉnh lại tất cả các mô hình & 90\% & 25.50\% & Tùng \\ \hline
	\end{tabularx}
\end{table}